\section{Exact Gaussian examples for the penalty-tail test}
\label{app:exact-ridge-examples}
Two constructed cases illustrate the strict large-penalty test and the role
of ties. Both use two prototypes at coordinates $0$ and
$d=\sqrt{2\log 2}$, bandwidth $\sigma=1$, and query $x=0$. Their Gaussian
cross-response is $1/2$, and the normalized query vector is
$r=(1,1/2)^\tr$. Select a class-$a$ observation at zero and a class-$b$
observation at $d$ as the two prototypes. The query has a unique nearest
prototype with label $a$.

For the first case, use training coordinates $(0,d,0)$ with labels $(a,b,b)$.
Direct multiplication gives
\[
 A=\begin{pmatrix}9/4&3/2\\3/2&3/2\end{pmatrix},
 \qquad Z[:,b]-Z[:,a]=\begin{pmatrix}1/2\\1\end{pmatrix}.
\]
The limiting competing margin is
$r^\tr(Z[:,b]-Z[:,a])=1$. The exact fitted margin is
\[
 r^\tr\bigl(W_{1,\lambda}[:,b]-W_{1,\lambda}[:,a]\bigr)
 =\frac{\lambda}{\det(A+\lambda I)}>0
 \qquad(\lambda>0).
\]
Thus every penalty region with no finite upper bound contains a strict
disagreement at this width.

For the second case, place five class-$a$ observations and one class-$b$
observation at zero, and one class-$a$ observation and six class-$b$
observations at $d$. Retain the same prototype labels and coordinates.
Here
\[
 A=\begin{pmatrix}31/4&13/2\\13/2&17/2\end{pmatrix},
 \qquad Z[:,b]-Z[:,a]=\begin{pmatrix}-3/2\\3\end{pmatrix}.
\]
The limiting competing margin equals zero. At every finite nonnegative
penalty, the exact fitted margin satisfies
\[
 r^\tr\bigl(W_{1,\lambda}[:,b]-W_{1,\lambda}[:,a]\bigr)
 =\frac{-63/4}{\det(A+\lambda I)}<0.
\]
The nearest-prototype class therefore wins throughout the finite penalty
path despite the exact limiting tie. Both Gram matrices are positive
definite, so their displayed denominators are positive for every
$\lambda\geq0$. These algebraic examples are independent of the image
evaluation and use only duplicates permitted by the model assumptions.
